\section{Introduction}
\label{submission:sec:intro}

The modern deep learning paradigm revolves around the pretraining-then-finetuning workflow. Large foundation models are pretrained on massive datasets to capture generalizable representations, and then specialized into task-specific experts via fine-tuning on downstream tasks~\cite{radford2021learning, devlin2018bert}. While fine-tuning produces highly accurate individual experts, deploying separate models for every task is prohibitively expensive in terms of storage, memory, and serving overhead, even when utilizing parameter-efficient transfer learning (PEFT) and adapter-based pipelines~\cite{hu21lora, dettmers2024qlora, houlsby2019parameter}. Consequently, multi-task learning (MTL) is highly desirable to consolidate diverse capabilities into a single neural network~\cite{caruana1997multitask, ruder2017overview}, or employing sparsely-gated Mixture-of-Experts (MoE) architectures~\cite{fedus2022switch, shazeer2017outrageously, jacobs1991adaptive}. 

However, joint multi-task training from scratch is often computationally restrictive and suffers from catastrophic forgetting or data privacy constraints. To circumvent this, \textit{model merging} has emerged as an attractive alternative~\cite{wortsman2022model, mcmahan2017communication, matena2022merging, ainsworth2023git, stoica2023zipit}. Model merging aims to combine several independently fine-tuned expert models sharing a common pretrained initialization into a single model, without requiring the original training data or additional fine-tuning.

Among merging paradigms, \textbf{Task Arithmetic}~\cite{ilharco2022editing} is widely used due to its conceptual simplicity. It computes the weight differences between each fine-tuned expert and the pretrained base model to form "task vectors," and then simply averages or sums these vectors before adding them back to the pretrained base. Despite its utility, standard linear parameter averaging is notoriously prone to failure because of the \textit{representation scale mismatch} across different layers and tasks. Because different experts are fine-tuned using heterogeneous schedules, learning rates, or data distributions, their parameter updates adapt at vastly different scales (i.e., divergent standard deviations and Frobenius norms). When averaged, the task with the largest weight standard deviation or the layer with the most aggressive updates disproportionately dominates the merged representations, causing severe interference and catastrophic performance degradation on the remaining tasks.

To combat this task interference, recent literature has seen a trend toward escalating complexity:
\begin{itemize}
    \item \textbf{SyMerge}~\cite{symerge2025synergistic} proposes test-time single-layer optimization using test-time gradient descent guided by expert teacher pseudo-labels, introducing optimization instability and high inference latency.
    \item \textbf{OrthoMerge}~\cite{orthomerge25manf} maps weight dynamics onto the Riemannian manifold of the orthogonal group using singular value decomposition (SVD), Cayley transformations, and Procrustes alignments.
    \item \textbf{SAIM}~\cite{saim25flat} couples momentum-based Sharpness-Aware Block Coordinate Descent (SA-BCD) during fine-tuning with Adaptive Isotropic Merging during the merging stage, which applies SVD to every layer weight tensor to balance the singular value spectrum.
\end{itemize}

From a \textit{minimalist} perspective, these approaches violate Occam's Razor. They introduce immense computational overhead, sensitive hyperparameters, complex multi-stage pipelines, or rely on active gradient descent and unlabeled data at test-time. We argue that model merging should remain highly lightweight, training-free, and easy to deploy. If a highly complex method's performance can be matched or exceeded by a simpler one, the simpler one is fundamentally superior.

\begin{figure}[t]
    \centering
    \includegraphics[width=0.98\linewidth]{results/fig1.png}
    \caption{\textbf{Multi-Task Merging Performance Comparison} on MNIST, FashionMNIST, and KMNIST image classification under heterogeneous training schedules. Our training-free, non-parametric \textbf{RMS-Scale} (73.22 $\pm$ 2.15\%) and \textbf{SD-Scale} (73.23 $\pm$ 2.19\%) perform outstandingly, slightly exceeding standard Task Arithmetic (72.50 $\pm$ 1.17\%) and Ties-Merging (71.77 $\pm$ 2.06\%) on average, while matching or exceeding SVD-based Isotropic Merging (73.13 $\pm$ 2.49\%) in strictly linear $O(N)$ time with just two lines of code and zero complex heuristics.}
    \label{fig:comparison_main}
\end{figure}

In this paper, we propose \textbf{Standard-Deviation Scaling (SD-Scale)} and its mathematically stable, non-translation-invariant counterpart \textbf{Root-Mean-Square Scaling (RMS-Scale)}. While standard deviation scaling is intuitive, standard deviation is translation-invariant because it subtracts the mean update. On low-variance parameter tensors, such as small biases, subtracting the mean can cause standard deviation to fall near zero, leading to division-by-zero or numerical instability when normalized. Although this vulnerability is largely theoretical and heavily diluted in practice—since bias parameters account for a negligible fraction of modern models—we introduce RMS-Scale as a mathematically stable, non-translation-invariant alternative. More practically, we also recommend simple weight-only scaling, which completely bypasses bias tensors while maintaining peak performance.

Our core insight is that task interference can be mitigated at the layer level through two simple scaling steps:
\begin{enumerate}
    \item \textbf{Layer-wise Normalization:} For every weight tensor in the model, we normalize each task vector to have unit Root-Mean-Square (RMS). This strips away magnitude disparities across tasks, ensuring that their updates contribute equally as isotropic direction vectors.
    \item \textbf{Scale Calibration:} We rescale the averaged normalized task vector by the average of the original RMS scales of the individual task vectors at that layer. This projects the balanced direction back onto the natural adaptation scale of the network, preventing representation distortion and retaining the appropriate layer-wise capacity.
\end{enumerate}

While standard RMS-Scale uses a global post-hoc scaling coefficient $\lambda$ tuned via a validation grid search, we introduce \textbf{Parameter-Free RMS-Scale (PF-RMS)} as our primary, completely parameter-free variant. When task-specific updates are averaged, parameter conflicts and partial task orthogonality cause a natural shrinkage of the merged task vector. PF-RMS analytically resolves this shrinkage at each layer by dynamically inverting the layer-wise alignment ratio $\alpha^l = \text{RMS}(\bar{\tau}_{\text{norm}}^l)$. This completely eliminates any post-hoc grid search, delivering robust, highly competitive multi-task capability with absolutely zero validation-set tuning.

Our main contributions are summarized as follows:
\begin{itemize}
    \item We demonstrate that standard task vector merging suffers heavily from representation scale mismatches across tasks and layers, leading to dominant task interference.
    \item We present \textbf{RMS-Scale} and its completely training-free and parameter-free variant \textbf{PF-RMS}. PF-RMS resolves scale mismatches and task conflicts through layer-wise RMS normalization and analytical scale calibration, completely bypassing any validation-set hyperparameter tuning.
    \item We evaluate our proposed methods on a rigorous multi-task image classification benchmark (MNIST, FashionMNIST, and KMNIST) fine-tuned under uncoordinated downstream schedules across 3 independent seeds. Our proposed SD-Scale (73.23\%) and RMS-Scale (73.22\%) slightly exceed standard Task Arithmetic (72.50\%) and Ties-Merging (71.77\%) on average, while our completely parameter-free \textbf{PF-RMS} achieves an outstanding \textbf{72.23\%} average accuracy completely out-of-the-box without requiring a validation set or any tuning, outperforming un-tuned Task Arithmetic (71.68\%) and un-tuned Ties-Merging (71.81\%).
    \item We mathematically and empirically demonstrate that RMS-Scale is equivalent to Frobenius-norm normalization scaled by parameter count, and show its superior numerical stability over standard deviation scaling.
    \item Through extensive ablation studies, we demonstrate the essential role played by both normalization and calibration, validating that minimalist scale calibration is robust, highly efficient, and sufficient.
\end{itemize}
